\label{bkm:Ref196027596}\index{BPatch_snippet: : : : : : :!}{

A snippet is an abstract representation of code to insert into a program. Snippets are defined by creating a new

instance of the correct subclass of a snippet. For example, to create a snippet to call a function, create a new

instance of the class BPatch_funcCallExpr. Creating a snippet does not result in code being inserted into an

application. Code is generated when a request is made to insert a snippet at a specific point in a program.

Sub-snippets may be shared by different snippets (i.e, a handle to a snippet may be passed as an argument to create two

different snippets), but whether the generated code is shared (or replicated) between two snippets is implementation

dependent.}





\bigskip



{

BPatch_type *getType\index{getType: : : : : : :!}()}



{

Return the type of the snippet. The BPatch_type system is described in section \ref{bkm:Ref161040683}.}





\bigskip



{

float getCost\index{getCost: : : : : : :!}()}



{

Returns an estimate of the number of seconds it would take to execute the snippet. The problems with accurately

estimating the cost of executing code are numerous and out of the scope of this document[2]. It is important to realize

that the returned cost value is, at best, an estimate.}



{

The rest of the classes are derived classes of the class BPatch_snippet.



{

BPatch_actualAddressExpr\index{BPatch_retExpr: : : : : : :!}()}



{

This snippet results in an expression that evaluates to the actual address of the instrumentation. To access the

original address where instrumentation was inserted, use BPatch_originalAddressExpr. Note that this actual address is

highly dependent on a number of internal variables and has no relation to the original address.}





\bigskip



{

BPatch_arithExpr\index{BPatch_arithExpr: : : : : : :!}(BPatch_binOp op, const BPatch_snippet &lOperand,\newline

const BPatch_snippet &rOperand)}



{

Perform the required binary operation. The available binary operators are:}



\begin{center}

\tablefirsthead{}

\tablehead{}

\tabletail{}

\tablelasttail{}

\begin{supertabular}{|m{1.1198599in}|m{4.3761597in}|}

\hline

{ Operator} &

{ Description}\\\hline

{ BPatch_assign} &

{ assign the value of rOperand to lOperand\\\hline

{ BPatch_plus} &

{ add lOperand and rOperand\\\hline

{ BPatch_minus} &

{ subtract rOperand from lOperand\\\hline

{ BPatch_divide} &

{ divide rOperand by lOperand}\\\hline

{ BPatch_times} &

{ multiply rOperand by lOperand}\\\hline

{ BPatch_ref} &

{ Array reference of the form lOperand[rOperand]}\\\hline

{ BPatch_seq} &

{ Define a sequence of two expressions (similar to comma in C)\\\hline

\end{supertabular}

\end{center}



\bigskip



{

BPatch_arithExpr(BPatch_unOp, const BPatch_snippet &operand)}



{

Define a snippet consisting of a unary operator. The unary operators are:}



\begin{center}

\tablefirsthead{}

\tablehead{}

\tabletail{}

\tablelasttail{}

\begin{supertabular}{|m{1.1198599in}|m{4.3761597in}|}

\hline

{ Operator} &

{ Description}\\\hline

{ BPatch_negate} &

{ Returns the negation of an integer}\\\hline

{ BPatch_addr} &

{ Returns a pointer to a BPatch_variableExpr}\\\hline

{ BPatch_deref} &

{ Dereferences a pointer\\\hline

\end{supertabular}

\end{center}



\bigskip



{

BPatch_boolExpr\index{BPatch_boolExpr: : : : : : :!}(BPatch_relOp op, const BPatch_snippet &lOperand,\newline

const BPatch_snippet &rOperand)}



{

Define a relational snippet. The available operators are:



\begin{center}

\tablefirsthead{}

\tablehead{}

\tabletail{}

\tablelasttail{}

\begin{supertabular}{|m{0.8788598in}|m{3.1087599in}|}

\hline

{ Operator} &

{ Function\\\hline

{ BPatch_lt} &

{ Return lOperand < rOperand\\\hline

{ BPatch_eq} &

{ Return lOperand == rOperand}\\\hline

{ BPatch_gt} &

{ Return lOperand > rOperand\\\hline

{ BPatch_le} &

{ Return lOperand <= rOperand\\\hline

{ BPatch_ne} &

{ Return lOperand != rOperand\\\hline

{ BPatch_ge} &

{ Return lOperand >= rOperand\\\hline

{ BPatch_and} &

{ Return lOperand && rOperand (Boolean and)\\\hline

{ BPatch_or} &

{ Return lOperand {\textbar}{\textbar} rOperand (Boolean or)\\\hline

\end{supertabular}

\end{center}

{

The type of the returned snippet is boolean, and the operands are type checked.}



{

BPatch_breakPointExpr\index{BPatch_breakPointExpr: : : : : : :!}()}



{

Define a snippet that stops a process when executed by it. The stop can be detected using the isStopped member

function of BPatch_process, and the program's execution can be resumed by calling the

continueExecution member function of BPatch_process.}





\bigskip



{

BPatch_bytesAccessedExpr\index{BPatch_bytesAccessedExpr : : : : : : :!}()}



{

This expression returns the number of bytes accessed by a memory operation. For most load/store architecture machines

it is a constant expression returning the number of bytes for the particular style of load or store. This snippet is

only valid at a memory operation instrumentation point.}





\bigskip



{

BPatch_constExpr\index{BPatch_constExpr: : : : : : :!}(signed int value)}



{

BPatch_constExpr\index{BPatch_constExpr: : : : : : :!}(unsigned int value)



{

BPatch_constExpr(signed long value)}



{

BPatch_constExpr(unsigned long value)}



{

BPatch_constExpr(const char *value)}



{

BPatch_constExpr(const void *value)}



{

BPatch_constExpr(long long value)}



{

Define a constant snippet of the appropriate type. The char* form of the constructor creates a constant string; the

null-terminated string beginning at the location pointed to by the parameter is copied into the

application's address space, and the BPatch_constExpr that is created refers to the location to

which the string was copied.}





\bigskip



{

BPatch_dynamicTargetExpr\index{BPatch_retExpr: : : : : : :!}()}



{

This snippet calculates the target of a control flow instruction with a dynamically determined target. It can handle

dynamic calls, jumps, and return statements.





\bigskip



{

BPatch_effectiveAddressExpr\index{BPatch_effectiveAddressesExpr : : : : : : :!}()}



{

Define an expression that contains the effective address of a memory operation. For a multi-word memory operation

(i.e. more than the “natural” operation size of the machine), the effective address is the base address of the

operation.





\bigskip



{

BPatch_funcCallExpr\index{BPatch_funcCallExpr: : : : : : :!}(const BPatch_function& func,\newline

const std::vector<BPatch_snippet*> &args)}



{

Define a call to a function. The passed function must be valid for the current code region. Args is a list of

arguments to pass to the function; the maximum number of arguments varies by platform and is summarized below. If type

checking is enabled, the types of the passed arguments are checked against the function to be called. Availability of

type checking depends on the source language of the application and program being compiled for debugging.}



\begin{center}

\tablefirsthead{}

\tablehead{}

\tabletail{}

\tablelasttail{}

\begin{supertabular}{|m{1.0462599in}|m{2.94136in}|}

\hline

{ Platform} &

{ Maximum number of arguments}\\\hline

{ AMD64/EMT-64} &

{ No limit}\\\hline

{ IA-32} &

{ No limit}\\\hline

{ POWER} &

{ 8 arguments}\\\hline

\end{supertabular}

\end{center}



\bigskip



{

BPatch_funcJumpExpr\index{funcJumpExpr: : : : : : :!} (const BPatch_function &func)}



{

This snippet has been removed; use BPatch_addressSpace::wrapFunction instead.}





\bigskip



{

BPatch_ifExpr\index{BPatch_ifExpr: : : : : : :!}(const BPatch_boolExpr &conditional,\newline

const BPatch_snippet &tClause,\newline

const BPatch_snippet &fClause)}



{

BPatch_ifExpr(const BPatch_boolExpr &conditional,\newline

const BPatch_snippet &tClause)}



{

This constructor creates an if statement. The first argument, conditional, should be a Boolean expression that will be

evaluated to decide which clause should be executed. The second argument, tClause, is the snippet to execute if the

conditional evaluates to true. The third argument, fClause, is the snippet to execute if the conditional evaluates to

false. This third argument is optional. Else-if statements, can be constructed by making the fClause of an if statement

another if statement.}





\bigskip



{

BPatch_nullExpr\index{Bpatch_nullExpr: : : : : : :!}()}



{

Define a null snippet. This snippet contains no executable statements.}





\bigskip



{

BPatch_originalAddressExpr\index{BPatch_retExpr: : : : : : :!}()}



{

This snippet results in an expression that evaluates to the original address of the point where the snippet was

inserted. To access the actual address where instrumentation is executed, use BPatch_actualAddressExpr.}





\bigskip



{

BPatch_paramExpr\index{Bpatch_paramExpr: : : : : : :!}(int paramNum)}



{

This constructor creates an expression whose value is a parameter being passed to a function. ParamNum specifies the

number of the parameter to return, starting at 0. Since the contents of parameters may change during subroutine

execution, this snippet type is only valid at points that are entries to subroutines, or when inserted at a call point

with the when parameter set to BPatch_callBefore.}





\bigskip



{

BPatch_registerExpr\index{BPatch_retExpr: : : : : : :!}(BPatch_register reg)}



{

BPatch_registerExpr(Dyninst::MachRegister reg)}



{

This snippet results in an expression whose value is the value in the register at the point of instrumentation.}





\bigskip



{

BPatch_retExpr\index{BPatch_retExpr: : : : : : :!}()}



{

This snippet results in an expression that evaluates to the return value of a subroutine. This snippet type is only

valid at BPatch_exit points, or at a call point with the when parameter set to BPatch_callAfter.}





\bigskip





\bigskip





\bigskip





\bigskip



{

BPatch_scrambleRegistersExpr\index{Bpatch_nullExpr: : : : : : :!}()}



{

This snippet sets all General Purpose Registers to the flag value.}





\bigskip



{

BPatch_sequence\index{BPatch_sequence: : : : : : :!}(const std::vector<BPatch_snippet*>

&items)}



{

Define a sequence of snippets. The passed snippets will be executed in the order in which they appear in items.





\bigskip



{

BPatch_shadowExpr\index{Bpatch_nullExpr: : : : : : :!}(bool entry,}



{

const BPatchStopThreadCallback &cb,}



{

const BPatch_snippet &calculation,}



{

bool useCache = false,}



{

BPatch_stInterpret interp = BPatch_noInterp)}



{

This snippet creates a shadow copy of the snippet BPatch_stopThreadExpr.}





\bigskip



{

BPatch_stopThreadExpr\index{Bpatch_nullExpr: : : : : : :!}(const BPatchStopThreadCallback &cb,}



{

const BPatch_snippet &calculation,}



{

bool useCache = false,}



{

BPatch_stInterpret interp = BPatch_noInterp)}



{

This snippet stops the thread that executes it. It evaluates a calculation snippet and triggers a callback to the user

program with the result of the calculation and a pointer to the BPatch_point at which the snippet was inserted.}





\bigskip



{

BPatch_threadIndexExpr()}



{

This snippet returns an integer expression that contains the thread index of the thread that is executing this snippet.

The thread index is the same value that is returned on the mutator side by BPatch_thread::getBPatchID.}





\bigskip



{

BPatch_tidExpr\index{BPatch_tidExpr: : : : : : :!}(BPatch_process *proc)}



{

This snippet results in an integer expression that contains the tid of the thread that is executing this snippet. This

can be used to record the threadId, or to filter instrumentation so that it only executes for a specific thread.}



{

BPatch_variableExpr\index{Bpatch_nullExpr: : : : : : :!}(char *in_name,}



{

BPatch_addressSpace *in_addSpace,}



{

AddressSpace *as,}



{

AstNodePtr ast_wrapper_,}



{

BPatch_type *type, void* in_address)}



{

BPatch_variableExpr\index{Bpatch_nullExpr: : : : : : :!}(BPatch_addressSpace *in_addSpace,}



{

AddressSpace *as,}



{

void *in_address,}



{

int in_register,}



{

BPatch_type *type,}



{

BPatch_storagestorage = BPatch_storageAddr,}



{

BPatch_point *scp = NULL)}



{

BPatch_variableExpr\index{Bpatch_nullExpr: : : : : : :!}(BPatch_addressSpace *in_addSpace,}



{

AddressSpace *as,}



{

BPatch_localVar *lv,}



{

BPatch_type *type,}



{

BPatch_point *scp)}



{

BPatch_variableExpr\index{Bpatch_nullExpr: : : : : : :!}(BPatch_addressSpace *in_addSpace,}



{

AddressSpace *ll_addSpace,}



{

int_variable *iv,}



{

BPatch_type *type)}



{

Define a variable snippet of the appropriate type. The first constructor is used to get function pointers; the second

is used to get forked copies of variable expression, used by malloc; the third is used for local variables; and the

last is used by BPatch_addressSpace::findOrCreateVariable().}



{

BPatch_whileExpr\index{Bpatch_nullExpr: : : : : : :!}(const BPatch_snippet &condition,}



{

const BPatch_snippet &body)



{

This constructor creates a while statement. The first argument, condition, should be a Boolean expression that will be

evaluated to decide whether body should be executed. The second argument, body, is the snippet to execute if the

condition evaluates to true.}





\bigskip



